\section{El concepto xG y su generalización}
\label{sec:xg}

\subsection{Origen y motivación}

El concepto \emph{Expected Goals} (xG) fue introducido por Sam Green en 2012 en un blog de Opta, formalizando una intuición previa: no todos los tiros tienen la misma probabilidad de convertirse en gol. Un mano-a-mano con el arquero a tres metros de la portería y un disparo de treinta metros con cinco defensores en medio son eventos cualitativamente distintos, aunque ambos cuenten como ``un tiro'' en las estadísticas tradicionales.

La idea es asignar a cada tiro una probabilidad estimada de gol, derivada empíricamente de cuántas veces tiros con características similares se convirtieron en gol a lo largo de un corpus histórico. Esa probabilidad es la xG del tiro. Sumando xG sobre los tiros de un partido o temporada se obtiene la métrica agregada $\mathrm{xG}$ de un jugador o equipo.

\subsection{Definición formal}

Sea $x_i$ un tiro y $\mathbf{c}_i \in \mathbb{R}^k$ el vector de contexto observado en el momento del tiro. Sea $Y_i \in \{0,1\}$ el indicador de gol ($Y_i = 1$ si el tiro fue gol). Un modelo xG estima:

\[
  \mathrm{xG}(x_i) \;=\; P(Y_i = 1 \mid \mathbf{c}_i)
\]

\noindent
en la práctica vía regresión logística, \emph{gradient boosting} o redes neuronales. Para un conjunto de tiros $\{x_1,\dots,x_n\}$ se define la xG agregada como:

\[
  \mathrm{xG}_{\mathrm{agg}} \;=\; \sum_{i=1}^{n} \mathrm{xG}(x_i)
\]

\noindent
y la diferencia $G - \mathrm{xG}_{\mathrm{agg}}$, donde $G = \sum_i Y_i$ es el número de goles realmente convertidos, se interpreta como \emph{calidad de definición} (\emph{finishing skill}): valores positivos indican que el jugador convierte más de lo que la jugada promedio sugiere; valores negativos, lo contrario.

\subsection{Features típicas del vector de contexto}

Las variables que componen $\mathbf{c}_i$ varían entre modelos, pero suelen incluir:

\begin{table}[h]
\centering
\begin{tabular}{ll}
\toprule
\textbf{Feature} & \textbf{Interpretación} \\
\midrule
Distancia a la portería          & Variable más importante \\
Ángulo de tiro                   & Apertura geométrica del marco \\
Parte del cuerpo                 & Pie hábil vs no hábil vs cabeza \\
Tipo de jugada previa            & Pase, centro, rebote, contraataque \\
Bajo presión defensiva           & Defensores entre tirador y arco \\
Posición del arquero             & Sólo en modelos con \emph{freeze-frame} \\
Número de defensores             & Sólo en modelos con \emph{freeze-frame} \\
Velocidad del balón              & Sólo con datos de tracking \\
\bottomrule
\end{tabular}
\caption{Features comunes en modelos xG. Las dos últimas requieren datos de tracking de jugadores que sólo se exponen comercialmente.}
\end{table}

El modelo xG de StatsBomb publicado en su \emph{open data} usa el \emph{freeze-frame}: una fotografía del momento del tiro con la posición de todos los jugadores, lo que permite calcular cuántos defensores están entre el tirador y la portería, dónde está el arquero, etc. Otros modelos (FBref, Understat) usan sólo la geometría del tiro y son por lo tanto menos precisos.

\subsection{Interpretación práctica}

\begin{table}[h]
\centering
\begin{tabular}{lc}
\toprule
\textbf{Tipo de tiro} & \textbf{xG aproximado} \\
\midrule
Penal & 0.76 \\
\emph{Tap-in} a dos metros de portería vacía & 0.95 \\
Cabezazo desde el corazón del área & 0.20 -- 0.40 \\
Cabezazo desde primer palo, ángulo cerrado & 0.05 \\
Disparo desde fuera del área ($\sim$25 m) & 0.03 -- 0.05 \\
Tiro libre directo desde 25 m & 0.06 \\
Mano a mano con el arquero & 0.40 -- 0.60 \\
\bottomrule
\end{tabular}
\caption{Rangos típicos de xG según el tipo de tiro. Calibrados con datos abiertos de StatsBomb.}
\end{table}

\subsection{Por qué xG predice mejor que goles}

Los goles son eventos raros y de varianza muy alta partido a partido. Para un mismo nivel real de habilidad, un delantero puede meter tres goles en un partido y ninguno en los siguientes cuatro. xG es una métrica de varianza mucho menor: la cantidad de chances de calidad generadas por un jugador es relativamente estable y refleja con más fidelidad su contribución sistemática al juego.

Empíricamente, en nuestros datos abiertos de torneos internacionales (FIFA World Cup 2018, 2022; UEFA Euro 2020, 2024; Copa América 2024; \emph{etcétera}, total 908 partidos), la correlación de Pearson entre goles y xG es $r \approx 0.62$, mientras que entre goles y tiros (sin ajustar por calidad) es $r \approx 0.44$. Ese delta de $0.18$ es la mejora que aporta el ajuste por calidad de chance: xG es un predictor de goles futuros más fuerte que los goles propios pasados.

\subsection{Generalización a cualquier evento contable}

La construcción anterior se generaliza directamente. Para cualquier evento discreto $M$ que se pueda observar y etiquetar binariamente (\emph{ocurre / no ocurre}), si existe un vector de contexto $\mathbf{c}$ predictivo, se puede definir:

\[
  \mathrm{x}\textsc{M}(x_i) \;=\; P(M_i = 1 \mid \mathbf{c}_i)
\]

\noindent
y su agregado correspondiente. Esto da pie al catálogo descrito en la Sección~\ref{sec:catalogo}.

Los principios de identificación de un buen candidato para una métrica \emph{expected} son tres:

\begin{enumerate}[label=\textbf{P\arabic*.}]
  \item \textbf{Granularidad por evento}: cada ocurrencia debe ser identificable individualmente (no sólo un \emph{count} agregado).
  \item \textbf{Contexto observable}: deben existir features pre-evento que predigan razonablemente la ocurrencia.
  \item \textbf{Aplicabilidad a predicción}: el evento debe ser repetible suficientes veces para que el agregado tenga señal.
\end{enumerate}
